% Unit 084 English translation of chapter6.tex, lines 2313--2457.
% Source: Methods of Algebra, Volume 2, commit
% 9a5803ff77dd3257484cb177f851a73770a59dd3 (CC BY 4.0).
% stable-unit-id: o014.aljabr2.chapter6.profinite-cohomology-a
% Coverage: Section 6.12 through smooth modules, the definition of profinite
% cohomology, and continuous standard complexes through its first theorem.

% segment-id: o014.li2.chapter6.profinite-cohomology-a.h001
\section{Cohomology of Profinite Groups}\label{sec:profinite-cohomology}
\hypertarget{o014.aljabr2.chapter6.profinite-cohomology-a}{ }

% segment-id: o014.li2.chapter6.profinite-cohomology-a.p001
This section presupposes the basic definitions and results on profinite groups
from \S\ref{sec:profinite-groups}.

% segment-id: o014.li2.chapter6.profinite-cohomology-a.p002
Let $G$ be a profinite group and $M$ a $G$-module. It is not difficult to
verify that the following conditions are equivalent:
\begin{enumerate}[(i)]
	\item $M=\bigcup_KM^K$, where $K$ ranges over the open normal subgroups of
	$G$;
	\item for every $x\in M$, the stabilizer subgroup $\Stab_G(x)$ is open;
	\item the action map $G\times M\to M$ is continuous when $M$ is given the
	discrete topology.
\end{enumerate}

% segment-id: o014.li2.chapter6.profinite-cohomology-a.d001
\begin{definition}
	\index{$G$-module!smooth}
	\index[sym1]{G-Mod-infty@$G\dcate{Mod}^\infty$}
	Let $G$ be a profinite group. If any of the conditions above holds, the
	$G$-module $M$ is called \emph{smooth}\footnote{In view of~(iii), many
	sources call it discrete.}. All smooth $G$-modules form a full subcategory
	$G\dcate{Mod}^\infty$ of $G\dcate{Mod}$.
\end{definition}

% segment-id: o014.li2.chapter6.profinite-cohomology-a.p003
Although the condition on $G$ itself is topological, characterization~(i) or
~(ii) of a smooth $G$-module is entirely algebraic: smoothness is a property
of a $G$-module, not additional structure. It follows readily that
$G\dcate{Mod}^\infty$ is an abelian subcategory of $G\dcate{Mod}$. The
operations of Definition~\ref{def:Res-Infl} have natural analogues here.

% segment-id: o014.li2.chapter6.profinite-cohomology-a.d002
\begin{definition-proposition}
	Let $\varphi:H\to G$ be a continuous homomorphism of profinite groups and
	let $M$ be a smooth $G$-module. Then the $H$-module $\varphi^*M$ is still
	smooth. This gives a functor
	\[ \varphi^*:G\dcate{Mod}^\infty\to H\dcate{Mod}^\infty. \]
\end{definition-proposition}
\begin{proof}
	For every open normal subgroup $K\lhd G$, its inverse image
	$\varphi^{-1}(K)\lhd H$ is again open and normal. If $x\in M^K$, then
	$x\in(\varphi^*M)^{\varphi^{-1}K}$. Thus
	\[ \varphi^*M=
	\bigcup_{L\lhd H:\,\text{open}}(\varphi^*M)^L, \]
	which proves that $\varphi^*M$ is smooth.
\end{proof}

% segment-id: o014.li2.chapter6.profinite-cohomology-a.p004
Note that the functor $\varphi^*$ is always exact. We single out two special
cases:
\begin{description}
	\item[Restriction] If $H$ is a closed subgroup of $G$, take $\varphi$ to
	be the inclusion homomorphism $H\hookrightarrow G$. The corresponding
	functor is denoted by
	\[ \Res^G_H:G\dcate{Mod}^\infty\to H\dcate{Mod}^\infty. \]
	\item[Inflation] For a closed normal subgroup $H\lhd G$, take $\varphi$ to
	be the quotient homomorphism $G\twoheadrightarrow G/H$. The corresponding
	functor is denoted by
	\[ \mathrm{Infl}^G_{G/H}:G/H\dcate{Mod}^\infty
	\to G\dcate{Mod}^\infty. \]

	The groups $H$ and $G/H$ mentioned above are automatically profinite, so
	these definitions make sense.
\end{description}
\index[sym1]{ResGH}
\index[sym1]{InflGH}

% segment-id: o014.li2.chapter6.profinite-cohomology-a.d003
\begin{definition-proposition}\label{def:smooth-part}
	For every $G$-module $M$, define
	$M^\infty:=\bigcup_KM^K$, where $K$ ranges over the open normal subgroups
	of $G$. This is a smooth $G$-module, and $M$ is smooth if and only if
	$M=M^\infty$.

	For every smooth $G$-module $M'$, we have
	\[ \Hom_G(M',M)=\Hom_G(M',M^\infty). \]
	In other words,
	$\text{inclusion}:G\dcate{Mod}^\infty
	\leftrightarrows G\dcate{Mod}:(\cdot)^\infty$ is an adjoint pair.
\end{definition-proposition}
\begin{proof}
	Clear.
\end{proof}

% segment-id: o014.li2.chapter6.profinite-cohomology-a.p005
The elements of $M^\infty$ are also called the \emph{smooth elements} of
$M$. Since $(\cdot)^\infty$ is the right adjoint of an exact functor, it
preserves injective objects
(Proposition~\ref{prop:adjoint-injective-projective}).

% segment-id: o014.li2.chapter6.profinite-cohomology-a.t001
\begin{corollary}
	The abelian category $G\dcate{Mod}^\infty$ has enough injective objects.
\end{corollary}
\begin{proof}
	Let $M$ be a smooth $G$-module and embed it in an injective $G$-module
	$I$. Then $M=M^\infty\hookrightarrow I^\infty$, while $I^\infty$ is an
	injective smooth $G$-module.
\end{proof}

% segment-id: o014.li2.chapter6.profinite-cohomology-a.r001
\begin{remark}
	In general, $G\dcate{Mod}^\infty$ does not have enough projective objects.
	This is the principal reason why this section studies cohomology rather
	than homology.
\end{remark}

% segment-id: o014.li2.chapter6.profinite-cohomology-a.p006
Since $G\dcate{Mod}^\infty$ has enough injective objects, we can discuss the
right derived functors of the invariants functor $(\cdot)^G$. Our approach is
first to give a direct definition that reduces to the case of finite groups,
and then to identify it with the right derived functors in
Proposition~\ref{prop:profinite-cohomology-identification}. For now we use
only classical language and do not invoke derived categories.

% segment-id: o014.li2.chapter6.profinite-cohomology-a.p007
First take a profinite group $G$ and an open normal subgroup $K$. For every
smooth $G$-module $M$, the submodule of $K$-invariants
$M^K=(\Res^G_KM)^K$ naturally becomes a $G/K$-module. This gives a functor
$(\cdot)^K:G\dcate{Mod}^\infty\to G/K\dcate{Mod}$ and the evident adjoint
pair
\begin{equation}\label{eqn:invariant-adjunction-profinite}
	\begin{tikzcd}[row sep=tiny]
		\mathrm{Infl}^G_{G/K}:G/K\dcate{Mod} \arrow[shift left,r] &
		G\dcate{Mod}^\infty:(\cdot)^K \arrow[shift left,l]\\
		\Hom_G(\mathrm{Infl}^G_{G/K}(M'),M) \arrow[equal,r] &
		\Hom_{G/K}(M',M^K).
	\end{tikzcd}
\end{equation}

% segment-id: o014.li2.chapter6.profinite-cohomology-a.p008
If open normal subgroups $K,K'$ of $G$ satisfy $K'\supset K$, then the
inclusion $M^{K'}\hookrightarrow M^K$ is equivariant with respect to the
oppositely directed quotient homomorphism $G/K'\twoheadleftarrow G/K$
(Definition~\ref{def:group-equivariance}). Hence for every
$n\in\Z_{\geq0}$ there is a canonical homomorphism
\[ \rho^{K'}_K:\Hm^n(G/K',M^{K'})\to\Hm^n(G/K,M^K). \]
If $K_1\subset K_2\subset K_3$ are open normal subgroups of $G$, then
$\rho^{K_2}_{K_1}\rho^{K_3}_{K_2}=\rho^{K_3}_{K_1}$. Thus the following
definition makes sense.

% segment-id: o014.li2.chapter6.profinite-cohomology-a.d004
\begin{definition}
	\index{group cohomology!profinite case}
	\index[sym1]{HnGM}
	Let $G$ be a profinite group and $M$ a smooth $G$-module. For every
	$n\in\Z_{\geq0}$, define
	\[ \Hm^n(G,M):=\varinjlim_K\Hm^n(G/K,M^K), \]
	where $K$ ranges over the open normal subgroups of $G$, and $\varinjlim$
	is taken with respect to the partial order
	$K'\preceq K\iff K\subset K'$ and the morphisms $\rho^{K'}_K$.
\end{definition}

% segment-id: o014.li2.chapter6.profinite-cohomology-a.p009
All open normal subgroups form a filtered partially ordered set: they are
closed under finite intersections. If $G$ is finite, then $\{1\}$ is its
unique maximal element, and $\Hm^n(G,M)$ reduces to the original version in
\S\ref{sec:group-coh-sub}. In general,
\[ \Hm^0(G,M)=\varinjlim_K(M^K)^{G/K}=M^G. \]

% segment-id: o014.li2.chapter6.profinite-cohomology-a.p010
Let $f:M_1\to M_2$ be a homomorphism of smooth $G$-modules. For every open
normal subgroup $K$, it restricts to
$f_K:M_1^K\to M_2^K$. Clearly, $K'\supset K$ implies
$\Hm^n(f_K)\rho^{K'}_K=\rho^{K'}_K\Hm^n(f_{K'})$, so for every
$n\in\Z_{\geq0}$ there is a homomorphism
\[ \Hm^n(f):\Hm^n(G,M_1)\to\Hm^n(G,M_2). \]

% segment-id: o014.li2.chapter6.profinite-cohomology-a.e001
\begin{example}
	Let $E|F$ be a Galois extension of fields. Since
	$E=\bigcup_{L|F}E^{\Gal(E|L)}=\bigcup_{L|F}L$, where $L|F$ ranges over
	the finite Galois subextensions, the action of $\Gal(E|F)$ on $E$ is
	smooth. Applying Example~\ref{eg:additive-90-finite} to every $L|F$
	gives
	\[ n\geq1\implies\Hm^n(\Gal(E|F),E)
	=\varinjlim_{L|F}\Hm^n(\Gal(L|F),L)=0. \]
\end{example}

% segment-id: o014.li2.chapter6.profinite-cohomology-a.p011
Just as in Proposition~\ref{prop:groupcoh-cochain} and
Remark~\ref{rem:nor-cochain} for groups without topology, the cohomology of a
profinite group can be computed by the standard complex or its normalized
version; the difference lies in the continuity condition.

% segment-id: o014.li2.chapter6.profinite-cohomology-a.d005
\begin{definition}
	\index{standard complex}
	\index[sym1]{CGM}
	For a smooth $G$-module $M$, give $M$ the discrete topology and define its
	\emph{standard complex} $C(G,M)=(C^n(G,M))_n$ by
	\[ C^n(G,M):=\{\text{continuous maps }f:G^n\to M\},
	\qquad n\in\Z_{\geq0}. \]
	Terms of negative degree are defined to be $0$, while
	$d^n:C^n(G,M)\to C^{n+1}(G,M)$ is given by
	\begin{multline*}
		(d^nf)(g_1,\ldots,g_{n+1})
		=g_1f(g_2,\ldots,g_{n+1})\\
		+\sum_{k=1}^n(-1)^kf(\ldots,g_kg_{k+1},\ldots)
		+(-1)^{n+1}f(g_1,\ldots,g_n).
	\end{multline*}

	\index[sym1]{CGMbar}
	Define the \emph{normalized standard complex} to be the following
	subcomplex $\overline{C}(G,M)$ of $C(G,M)$
	(see Remark~\ref{rem:nor-cochain}):
	\[ \overline{C}^n(G,M)=\left\{\begin{array}{r|l}
		f\in C^n(G,M)&\exists 1\leq k\leq n,\; g_k=1_G\\
		&\implies f(g_1,\ldots,g_n)=0
	\end{array}\right\}. \]
\end{definition}

% segment-id: o014.li2.chapter6.profinite-cohomology-a.p012
For every open normal subgroup $K\lhd G$, there are evident inclusions of
complexes
$C(G/K,M^K)\hookrightarrow C(G,M)$ and
$\overline{C}(G/K,M^K)\hookrightarrow\overline{C}(G,M)$.

% segment-id: o014.li2.chapter6.profinite-cohomology-a.p013
Since $M$ is discrete, a map $f:G^n\to M$ is continuous if and only if it is
locally constant. Compactness of $G$ then gives an open normal subgroup
$K_1\lhd G$ such that $f$ factors through $(G/K_1)^n\to M$. In particular,
$f$ assumes only finitely many values; smoothness gives an open normal
subgroup $K_2\lhd G$ such that the values of $f$ lie in $M^{K_2}$. Set
$K:=K_1\cap K_2$. It follows that
\[ C^n(G,M)\simeq\varinjlim_{K\lhd G\;\text{open}}C^n(G/K,M^K),\qquad
	\overline{C}^n(G,M)\simeq
	\varinjlim_{K\lhd G\;\text{open}}\overline{C}^n(G/K,M^K). \]
As $n$ varies, these give isomorphisms of complexes.

% segment-id: o014.li2.chapter6.profinite-cohomology-a.p014
As before, define the submodules of $C^n(G,M)$
\[ Z^n(G,M):=\Ker(d^n),\qquad B^n(G,M):=\Image(d^{n-1}), \]
with $B^0(G,M):=\{0\}$. Then
$\Hm^n(C(G,M))=Z^n(G,M)/B^n(G,M)$. Define the normalized versions
$\overline{Z}^n(G,M)$ and $\overline{B}^n(G,M)$ similarly.

% segment-id: o014.li2.chapter6.profinite-cohomology-a.t002
\begin{theorem}\label{prop:profinite-group-standard}
	Let $G$ be a profinite group. For all $n\in\Z_{\geq0}$ and every smooth
	$G$-module $M$, there are canonical isomorphisms
	\[ \Hm^n(C(G,M))\simeq\Hm^n(G,M)
	\simeq\Hm^n(\overline{C}(G,M)). \]
\end{theorem}
\begin{proof}
	Cohomology takes values in $\cate{Ab}$, or in $\Bbbk\dcate{Mod}$ if a
	coefficient ring $\Bbbk$ has been specified. Filtered colimits are exact
	in these categories \cite[Theorem 6.2.2]{Li1}. Hence
	\[ Z^n(G,M)=\varinjlim_KZ^n(G/K,M^K),\qquad
	B^n(G,M)=\varinjlim_KB^n(G/K,M^K), \]
	and therefore
	\begin{multline*}
		\Hm^n(C(G,M))
		=\dfrac{\varinjlim_KZ^n(G/K,M^K)}
		{\varinjlim_KB^n(G/K,M^K)}\simeq\\
		\varinjlim_K\frac{Z^n(G/K,M^K)}{B^n(G/K,M^K)}
		=\varinjlim_K\Hm^n(C(G/K,M^K)).
	\end{multline*}
	Apply Proposition~\ref{prop:groupcoh-cochain} to the last term to obtain
	$\varinjlim_K\Hm^n(G/K,M^K)$, namely $\Hm^n(G,M)$.
% Editorial correction carried from the Indonesian witness: the Chinese source
% has $\Hm^n(G,K)$ here, with $K$ in place of the coefficient module $M$.

	The case of $\overline{C}(G,M)$ is similar.
\end{proof}
